\documentclass[10pt,letterpaper]{article}
\usepackage{cornell-notes}

\newcommand{\CornellDocumentTitle}{CISSP CBK Cornell Notes: Domain 4 -- Communication and Network Security}
\title{\CornellDocumentTitle}
\author{Arthur Deane and Aaron Kraus}
\date{\today}

\setDocTitle{\CornellDocumentTitle}
\setDocSubtitle{Domain 4: Communication and Network Security}
\setDocVersion{Sixth Edition \textbullet\ 2021}
\setDocStatus{Study Companion}
\setDocOwner{Security Certification Study Collection}
\setDocAuthor{Arthur Deane and Aaron Kraus}
\setDocDate{\today}
\setDocSummary{Cornell-format study and review notes for Domain 4 of the CISSP CBK reference.}

\setCornellCollection{CISSP CBK Cornell Notes}
\setCornellUnitType{Domain}
\setCornellUnitNumber{4}
\setCornellUnitTitle{Communication and Network Security}
\setCornellSourceLabel{CBK}
\setCornellSourceTitle{THE OFFICIAL (ISC)\textsuperscript{2} CISSP CBK REFERENCE}
\setCornellSourceAuthor{Arthur Deane and Aaron Kraus}
\setCornellSourceEdition{Sixth Edition \textbullet\ 2021}
\setCornellSourceRange{pp. 283--374}
\setCornellCompanionLabel{Study and review companion}

\begin{document}
\maketitle

\section*{Domain Overview}
Domain 4 explains how data moves and how security architecture controls that movement. It combines layered network models, IP and secure protocols, segmentation, wired/wireless transmission, network security devices, endpoint/NAC controls, remote access, voice/collaboration, virtual networks, and third-party connectivity.

\begin{summarybox}{Review Objectives}
Map technologies and threats to OSI/TCP-IP layers; explain IP networking and secure protocols; design segmentation and wireless security; distinguish network devices/firewall types; secure media, endpoints, NAC, and mobile systems; and choose controls for voice, collaboration, remote access, VPNs, data communications, virtual networks, and external partners.
\end{summarybox}

\section{Secure Network Architecture and Layered Models}
\source{pp. 283--318}
\begin{cornell}
\noterow{Why use layered models?}{Layering separates communication functions into abstractions so protocols can interoperate and security problems can be localized. A control at one layer does not automatically protect other layers; understand where trust, addressing, routing, sessions, and applications actually operate.}
\noterow{OSI Layer 1 -- Physical}{Transmits raw bits over electrical/optical/radio media. Think cabling, connectors, signal strength, repeaters/hubs, interference, eavesdropping, jamming, physical access, and environmental protection.}
\noterow{OSI Layer 2 -- Data Link}{Moves frames on a local link using MAC addressing; Ethernet, switching, VLANs and link-layer wireless functions live here. Threats include MAC spoofing, ARP-related/local attacks, VLAN misconfiguration/hopping, and unauthorized attachment.}
\noterow{OSI Layer 3 -- Network}{Routes packets among networks using logical addressing. IP, ICMP, routers, and IPsec are associated here. Security focuses on route/address trust, filtering, segmentation, spoofing, and controlled inter-network flow.}
\noterow{OSI Layer 4 -- Transport}{Provides end-to-end transport, ports, segmentation/reassembly, and (with TCP) connection reliability/state; UDP is connectionless. Firewalls often filter ports/state here. Flooding, scanning, session manipulation, and exposed services are key concerns.}
\noterow{OSI Layer 5 -- Session}{Establishes, manages, synchronizes, and terminates communication sessions. Session hijacking, weak authentication, and legacy session/protocol weaknesses are typical concerns.}
\noterow{OSI Layer 6 -- Presentation}{Represents/transforms data for applications, including encoding, serialization, compression, and conceptually encryption. Problems in parsers/format handling can cross trust boundaries.}
\noterow{OSI Layer 7 -- Application}{User-facing network services/protocol logic such as HTTP, SMTP, FTP, DNS/SNMP-related application functions. Security concerns include vulnerable services, malicious inputs, weak application authentication, and protocol misuse.}
\noterow{TCP/IP model}{Commonly viewed as \textbf{Link/Network Access}, \textbf{Internet}, \textbf{Transport}, and \textbf{Application}. It maps practical Internet protocols to fewer layers than OSI. Use OSI for conceptual troubleshooting and TCP/IP for actual protocol-stack reasoning.}
\noterow{IP addressing and routing}{IP provides logical addressing and routed delivery; routers forward packets between networks. Subnetting creates smaller address/routing domains; gateways/default routes move traffic beyond the local subnet. Security architecture uses these boundaries for policy enforcement.}
\noterow{IPv4 / IPv6}{IPv4 uses 32-bit addresses and often NAT due to address scarcity. IPv6 uses 128-bit addressing and changes neighbor discovery, configuration, and operational assumptions. Security teams must explicitly govern both stacks; unmanaged IPv6 can become an alternate path around IPv4-focused controls.}
\noterow{TCP vs. UDP}{TCP is connection-oriented and tracks state/reliability with a handshake; UDP is connectionless and lower-overhead. Choose controls that understand protocol semantics: stateful filtering helps TCP, while UDP security relies more on application/protocol validation and rate/flow controls.}
\noterow{Common infrastructure protocols}{DNS maps names, DHCP assigns network configuration, ARP/neighbor functions resolve local addresses, ICMP carries network-control/error information, and SNMP manages devices. Each is operationally important and can be abused if unauthenticated, exposed, or misconfigured.}
\noterow{Secure Shell (SSH)}{Encrypted replacement for Telnet for remote command-line access/tunneling. Use SSH-2; the chapter identifies SSH-1 as insecure. Protect host keys/private keys, authenticate endpoints, restrict forwarding, and log privileged use.}
\noterow{TLS}{Replaced insecure SSL for secure client-server communications. TLS supports encryption/integrity and certificate-based authentication; it protects higher-layer protocols such as HTTPS and can be used for VPN/tunneling designs. Disable obsolete SSL/TLS versions and weak cipher configurations.}
\noterow{Kerberos}{Ticket-based protocol that lets principals prove identity across an untrusted network without repeatedly exposing reusable credentials. Time synchronization and protection of the Key Distribution Center/ticket material are critical.}
\noterow{IPsec}{Network-layer security suite providing authentication/integrity and confidentiality for IP. \textbf{AH} authenticates/integrity-protects but does not encrypt payload content; \textbf{ESP} can provide confidentiality plus authentication/integrity. Security Associations define negotiated parameters.}
\noterow{IPsec transport vs. tunnel}{\textbf{Transport mode} protects the IP payload while retaining the original outer IP header. \textbf{Tunnel mode} encapsulates/protects the original IP packet inside a new packet, making it suitable for site-to-site VPNs/security gateways.}
\noterow{IKE}{Internet Key Exchange establishes authenticated security associations/session keys for IPsec, commonly using certificates and Diffie-Hellman-style key agreement. Prefer current secure configuration and avoid downgrade/weak shared-secret designs.}
\noterow{Multilayer protocol implications}{Encapsulation places one protocol inside another, and security inspection at an outer layer may miss malicious or unauthorized inner traffic. Tunneling, encryption, fragmentation, and translation can obscure payloads; enforcement points must understand the complete stack they are expected to control.}
\noterow{Converged protocols}{Voice, video, storage, control, and ordinary data increasingly share IP infrastructure. Convergence reduces separate infrastructure but creates shared failure paths, QoS competition, and broader attack impact. Segment and prioritize traffic based on security and availability needs.}
\noterow{Microsegmentation}{Create small policy zones around workloads/resources and re-authenticate/re-authorize cross-zone access. Software-defined enforcement can restrict east-west lateral movement and support zero-trust designs, but policy/identity inventory and centralized control-plane security are essential.}
\end{cornell}

\section{Wireless, Cellular, and Content Distribution}
\source{pp. 319--334}
\begin{cornell}
\noterow{Wireless threat model}{Radio signals extend beyond walls and can be monitored/jammed/spoofed without physical cable access. Control access point placement/power, use strong authentication/encryption, detect rogue/evil-twin APs, segment wireless clients, and monitor the RF/network environment.}
\noterow{802.11 family}{Wi-Fi is governed by IEEE 802.11 amendments. Generations change frequency bands and performance, but higher speed does not imply stronger security. Do not confuse \textbf{802.11} wireless networking with \textbf{802.1X} port-based access control.}
\noterow{WEP}{Legacy RC4-based wireless protection with static/shared key weaknesses; easily broken and unsuitable. Recognize it as a historical mechanism that should be replaced, not tuned.}
\noterow{WPA / TKIP}{Interim improvement over WEP using TKIP and per-packet keying, but now considered insecure/legacy. The chapter advises moving from WEP/WPA and, where possible, older deployments toward current protection.}
\noterow{WPA2 / 802.11i}{Introduced AES/CCMP and personal vs. enterprise modes. Enterprise deployments can use 802.1X/EAP with centralized authentication. WPA2 is stronger than WEP/WPA but has known attack considerations and has been superseded by WPA3.}
\noterow{WPA3}{Current generation in this edition, adding stronger/individualized protection and weak-password mitigation. Use current supported modes and strong enterprise authentication where appropriate; transition older clients rather than weakening the entire WLAN.}
\noterow{802.1X / EAP}{Port-based Network Access Control framework using a supplicant, authenticator (switch/AP), and authentication server. EAP supports multiple methods; certificate-based EAP-TLS provides strong mutual authentication when certificates are well managed.}
\noterow{Bluetooth / short-range}{Short range reduces but does not remove risk. Control discoverability/pairing, authentication, supported profiles/services, device patching, and coexistence/interference; sensitive use may require policy restrictions.}
\noterow{Cellular networks}{Geographically distributed wireless cells and carrier infrastructure enable mobility, but radio and provider dependency create interception, spoofing, roaming, device, and availability risks. Sensitive applications still require end-to-end application/VPN protections rather than assuming the carrier link is secure.}
\noterow{Content Delivery Network (CDN)}{Distributes/caches content closer to users for performance and resilience and can absorb some traffic attacks. Security requires origin protection, trusted configuration, TLS/certificate management, cache rules, logging, provider assurance, and awareness that a compromised CDN/control plane affects many users.}
\end{cornell}

\section{Secure Network Components}
\source{pp. 335--356}
\begin{cornell}
\noterow{Firewall purpose}{Mediates traffic across a network boundary according to policy/ACL rules. Firewalls need correct placement, configuration, patching, protected administration, change control, logging, and fail behavior. They do not eliminate threats that never cross the enforcement point.}
\noterow{Static packet filter}{Layer-3-oriented screening of packet attributes such as addresses/protocols. Fast and simple, but little context and no inherent authentication; spoofing and complex application behavior can evade simplistic rules.}
\noterow{Stateful inspection}{Tracks connection state across network/transport layers so decisions consider whether packets belong to expected sessions. More context than stateless filtering; state tables and configuration still need protection and capacity.}
\noterow{Application-level firewall}{Layer-7 proxy/inspection understands application protocols and can apply granular content/command rules. Deeper inspection improves context but adds processing, complexity, and possible blind spots for encrypted/unknown traffic.}
\noterow{Circuit-level firewall}{Session-layer gateway validates session establishment/connection behavior rather than deeply inspecting application payload. Can mask internal addressing while providing less content awareness.}
\noterow{NGFW and deployment}{Next-generation firewalls combine traditional filtering with application awareness and IDS/IPS-like functions. Multihomed/bastion/screened-subnet designs create mediated zones. Cloud/SDN/FWaaS designs virtualize enforcement but still require least privilege and policy governance.}
\noterow{Repeaters / amplifiers}{Layer 1 devices extend or regenerate physical signals; they provide essentially no traffic intelligence or segmentation. Physical security and media protection dominate.}
\noterow{Hub}{Layer 1 multiport repeater that repeats traffic to ports, creating a shared collision/listening domain. Little security capability; switches replace hubs in modern controlled networks.}
\noterow{Bridge}{Layer 2 device that connects/filters local network segments using data-link information. It can reduce unnecessary traffic compared with a repeater but is largely subsumed by modern switches.}
\noterow{Switch}{Primarily Layer 2 intelligent forwarding based on MAC addresses; creates separate collision domains and can implement VLANs, port security, 802.1X, monitoring, and other controls. Protect management and prevent VLAN/Layer-2 attacks.}
\noterow{Router}{Layer 3 device that forwards packets between IP networks/subnets. Routing and ACLs create important policy boundaries; protect route information, management access, configuration, and availability.}
\noterow{Gateway}{Connects dissimilar networks/protocols or acts as an intermediary at higher layers. Translation/mediation creates a useful control point but also a complex trusted component that must validate both sides.}
\noterow{Proxy}{Intermediary that receives requests and makes them on behalf of another party; can filter, cache, inspect, authenticate, and conceal internal structure. Forward proxies represent clients; reverse proxies represent/protect servers.}
\noterow{Transmission media}{Copper, fiber, and wireless have different interception, interference, distance, cost, and physical-security characteristics. Protect cable routes/closets and use encryption where disclosure across the medium is unacceptable; fiber is harder, not impossible, to tap.}
\noterow{Network Access Control (NAC)}{Decide whether users/devices may connect based on identity, authentication and possibly posture/compliance. 802.1X/EAP can enforce admission at wired/wireless ports; quarantine/remediation zones restrict noncompliant endpoints.}
\noterow{Endpoint security}{Layer controls on laptops/desktops/servers/IoT: secure configuration, patching, malware/EDR protection, host firewall, disk encryption, least privilege, application control, device control, logging, and centralized management. The endpoint remains a trust boundary even behind a firewall.}
\noterow{Mobile devices}{Mobility increases loss/theft, untrusted-network, app, sensor, cloud-sync, and mixed personal/business risks. Use device management, encryption, screen lock, strong authentication, remote lock/wipe, app restrictions, patching, containerization/work profiles, and data minimization.}
\end{cornell}

\begin{exambox}{Layer-to-Device Cue}
Hub/repeater = Layer 1; bridge/switch = Layer 2; router = Layer 3. Firewalls/proxies can operate at multiple layers depending on type. Always reason from what the device actually inspects, not its marketing label.
\end{exambox}

\section{Secure Communication Channels}
\source{pp. 357--374}
\begin{cornell}
\noterow{Voice / VoIP}{Voice moved from dedicated telephony to IP, inheriting data-network threats plus toll fraud, eavesdropping, spoofing, call disruption, and signaling abuse. Segment voice, protect call managers, authenticate endpoints, secure signaling/media, patch devices, and preserve emergency/QoS requirements.}
\noterow{Multimedia collaboration}{Email, chat, file sharing, conferencing, and presence tools carry sensitive content through complex cloud/endpoints. Apply identity/MFA, encryption, retention/DLP, external-sharing policy, malware/link controls, meeting access controls, and logging.}
\noterow{Remote-access principle}{Remote access extends organizational trust to devices/networks outside physical control. Authenticate strongly, encrypt transport, enforce device posture/least privilege, monitor sessions, restrict split tunneling/administration as policy requires, and be able to revoke access quickly.}
\noterow{Remote authentication protocols}{The chapter surveys PAP, CHAP, EAP/PEAP/LEAP, RADIUS, and TACACS+. Recognize weak legacy password exchanges; prefer protected/centralized authentication and certificate/MFA-based approaches where supported.}
\noterow{VPN}{Creates a protected logical connection across an untrusted network. VPN encryption is only one control: the remote endpoint, identity, authorization, route/split-tunnel policy, client configuration, and monitoring determine whether the tunnel becomes a safe extension or a path for compromise.}
\noterow{Tunneling}{Encapsulates one protocol inside another so traffic can traverse an intermediate network. Tunnels can provide routing or security but may also hide traffic from controls; terminate and inspect them at governed trust boundaries.}
\noterow{VPN protocol cue}{IPsec operates at the network layer and supports transport/tunnel modes; TLS-based VPNs protect sessions/application access; older PPTP-related designs are legacy/weak; L2TP provides tunneling and is commonly paired with IPsec for security.}
\noterow{Frame Relay}{Legacy WAN packet-switching technology using virtual circuits between DTE/DCE devices. Recognize the concept of logical circuits over shared provider infrastructure and the need for additional protection if confidentiality/integrity are required.}
\noterow{ATM}{Asynchronous Transfer Mode uses fixed-size cells and supports QoS classes for voice/data. It is a legacy high-speed WAN/backbone technology; security requirements remain independent of the transport's performance guarantees.}
\noterow{MPLS}{Multiprotocol Label Switching forwards traffic along label-switched paths and enables provider VPN/QoS engineering. Logical separation by a carrier is not equivalent to end-to-end encryption; apply additional cryptographic protection when the threat model requires it.}
\noterow{Virtualized networks}{Software-defined networking, virtual switches/routers/firewalls, overlays, and cloud networks detach topology from physical devices. Secure centralized controllers/APIs, management credentials, tenant segmentation, policy automation, and monitoring; configuration errors can propagate rapidly.}
\noterow{Third-party connectivity}{Treat partners/vendors as separate trust domains. Define contractual/security requirements; use dedicated, least-privilege connections; authenticate/encrypt; segment; restrict routes/ports; log/monitor; reassess periodically; and remove access promptly when the relationship or need ends.}
\end{cornell}

\section{Critical Comparisons}
\begin{center}\small
\begin{tabularx}{\textwidth}{>{\bfseries\color{Navy}\raggedright\arraybackslash}p{0.22\textwidth}>{\raggedright\arraybackslash}X>{\raggedright\arraybackslash}X}
\toprule Concept & First idea & Distinguishing idea\\\midrule
TCP / UDP & Stateful, reliable connection & Connectionless, lower overhead\\
Switch / router & Layer 2 local frame forwarding & Layer 3 forwarding between IP networks\\
WEP/WPA / WPA2/WPA3 & Legacy/insecure generations & AES/CCMP and newer stronger wireless protection\\
AH / ESP & Authentication/integrity & Encryption plus authentication/integrity capabilities\\
IPsec transport / tunnel & Protect payload & Encapsulate/protect original packet\\
Packet filter / stateful & Per-packet attributes & Tracks connection state\\
Application firewall / circuit & Deep application proxy/inspection & Session validation with less payload awareness\\
\bottomrule
\end{tabularx}
\end{center}

\section{Key Terms and Acronyms}
\begin{summarybox}{Rapid-Review Glossary}
\textbf{ACL} -- Access Control List \quad
\textbf{AH} -- Authentication Header \quad
\textbf{CDN} -- Content Delivery/Distribution Network \quad
\textbf{EAP} -- Extensible Authentication Protocol \quad
\textbf{ESP} -- Encapsulating Security Payload \quad
\textbf{IKE} -- Internet Key Exchange \quad
\textbf{IPsec} -- Internet Protocol Security \quad
\textbf{MPLS} -- Multiprotocol Label Switching \quad
\textbf{NAC} -- Network Access Control \quad
\textbf{OSI} -- Open Systems Interconnection \quad
\textbf{SSH} -- Secure Shell \quad
\textbf{TLS} -- Transport Layer Security \quad
\textbf{VPN} -- Virtual Private Network \quad
\textbf{WPA3} -- Wi-Fi Protected Access 3.
\end{summarybox}

\section{High-Yield Review}
\begin{itemize}
\item Map a security control to the layer and data it actually sees.
\item Segmentation and microsegmentation limit trust and lateral movement; they require policy enforcement between zones.
\item SSH replaces Telnet; TLS replaces SSL; current IPsec uses AH/ESP plus security associations and IKE.
\item Modern WLAN security favors WPA3 and strong 802.1X/EAP-based enterprise authentication.
\item Firewalls are boundary controls, not complete network security; internal and encrypted paths must also be governed.
\item Remote access and third-party links must combine strong identity, encrypted transport, least privilege, endpoint posture, segmentation, and monitoring.
\end{itemize}

\section{Active-Recall Questions}
\begin{enumerate}
\item Map all seven OSI layers from lowest to highest.
\item What is the key difference between a Layer 2 switch and Layer 3 router?
\item How do IPsec AH and ESP differ?
\item How do IPsec transport and tunnel modes differ?
\item Why can multilayer encapsulation defeat a poorly placed control?
\item Why is 802.1X not the same as 802.11?
\item Order WEP, WPA, WPA2, WPA3 from oldest to newest and identify the current preference in this edition.
\item Compare stateless, stateful, application-level, and circuit-level firewalls.
\item What additional controls are needed even when a VPN encrypts remote traffic?
\item Why is provider/MPLS separation not automatically equivalent to encryption?
\item What does NAC decide before admitting an endpoint?
\item What should happen to a third-party connection when the business need ends?
\end{enumerate}

\subsection*{Answer Key}
\begin{enumerate}
\item Physical, Data Link, Network, Transport, Session, Presentation, Application.
\item A switch forwards frames on local Layer-2 segments/VLANs; a router forwards IP packets between Layer-3 networks.
\item AH authenticates/integrity-protects; ESP adds confidentiality and can provide authentication/integrity.
\item Transport protects the payload; tunnel protects/encapsulates the original packet inside a new packet.
\item A control may see only the outer protocol while unauthorized/malicious content is hidden in an inner tunnel/encrypted payload.
\item 802.11 is the Wi-Fi family; 802.1X is port-based network access control/authentication.
\item WEP, WPA, WPA2, WPA3; prefer WPA3 where possible in this edition.
\item Stateless checks packet attributes; stateful tracks connections; application-level understands application traffic; circuit-level validates sessions.
\item Strong identity/MFA, endpoint posture, authorization/least privilege, route policy, client security, logging/monitoring, and revocation.
\item It provides logical routing separation but not necessarily cryptographic confidentiality/integrity against all provider/path threats.
\item Whether device/user identity and required posture/policy conditions authorize network admission or quarantine.
\item Revoke/remove it promptly and update routes, credentials, rules, inventories, and monitoring.
\end{enumerate}

\section{Cornell Summary}
\begin{summarybox}{Domain 4 Summary}
Communication security is the controlled movement of information across layered, often untrusted infrastructure. Models explain where functions live; segmentation creates enforceable trust boundaries; secure protocols protect channels; wireless and remote access require explicit authentication and encryption; network devices mediate traffic according to the depth they can inspect; and virtual/third-party networks extend those same principles beyond owned hardware. Design each path so identity, authorization, confidentiality, integrity, availability, and monitoring remain intact from endpoint to endpoint.
\end{summarybox}

\section{Final Domain Checklist}
\begin{itemize}[label=\(\square\)]
\item I can map OSI/TCP-IP layers and principal devices/protocols.
\item I can explain IP addressing/routing, TCP/UDP, and infrastructure protocols.
\item I can select SSH, TLS, Kerberos, IPsec, and IKE for the proper role.
\item I can explain microsegmentation and wireless generations/802.1X.
\item I can distinguish firewall types and Layer 1/2/3 devices, gateways, and proxies.
\item I can secure transmission media, NAC, endpoints, and mobile devices.
\item I can design secure voice/collaboration/remote-access/VPN communications.
\item I can manage virtualized and third-party connectivity as explicit trust boundaries.
\end{itemize}
\vfill
{\footnotesize\color{Teal}Prepared as a study companion to Deane \& Kraus, \textit{The Official (ISC)\textsuperscript{2} CISSP CBK Reference}, 6th ed. Content is paraphrased and synthesized for study.}

\end{document}
